\scanpage{41}
% L1
Nous allons expliciter plus bas la détermination de $\sigma^0(X)$ et $\sigma^1(X)$

% L3
Il y a lieu également de préciser les opérations $\vee$ et $\lambda^i$. Pour la première, notons que
\[
\deleted{\mathbb Z}\mathbb Q(-i)\,[M_i^+(X)]^\vee\subset M_i^+(X)
\]
d'où on conclut de proche en proche que $x\mapsto\overset{\vee}{x}\xi^i$ définit sur $\sigma_i$ une involution. Mais en fait il résultera de la théorie \unsure{projective} que \uline{cette dernière est l'identité}, i.e. \uline{l'opération} $\vee$ \uline{dans} \deleted{$\unsure{A}$} $M_i(X)$ \deleted{\unsure{est la dualité}} \uline{n'est autre que} $x\mapsto\xi^{-i}x$, \deleted{\unsure{d'où} $x\in\unsure{M_i}(X)\Longrightarrow\overset{\vee}{x}=\xi^{\unsure{-i|i}}x$}

% L14
Quant aux opérations $\lambda^i$, leur détermination devrait se faire en termes de groupes finis d'opérateurs (tels les $\unsure{\mathfrak S|G}_p$ \dots).
% archive sheet number: 40
